\documentclass[10pt,a4paper]{amsart}
\usepackage[margin=19mm]{geometry}
\usepackage{amsmath,amssymb,mathtools}
\usepackage[scr=rsfs,cal=euler]{mathalfa}
\usepackage{xfrac}
\usepackage[unicode,hidelinks,bookmarks=false]{hyperref}
\newcommand{\ie}{i.e.\ }
\newcommand{\cf}{cf.\ }
\newcommand{\resp}{respectively\ }
\newcommand{\Subsection}{\subsection}
\providecommand{\nobreakdash}{\nobreak\hbox{-}}
\setlength{\emergencystretch}{2em}
\multlinegap0pt
\usepackage{amssymb}
\usepackage{enumerate}
\usepackage{natbib}
\usepackage{float}
\usepackage{bbm}
\newtheorem{proposition}{Proposition}
\title{Fully nonlinear prescribed curvature problems on closed manifolds with negative curvature}
\author{Jiaogen Zhang}
\date{Source: arXiv:2510.06577v2 (CC BY 4.0)}
\begin{document}
\maketitle

\noindent\emph{Source and licence.} Fully nonlinear prescribed curvature problems on closed manifolds with negative curvature. Version arXiv:2510.06577v2 (CC BY 4.0), distributed by arXiv under the Creative Commons Attribution 4.0 International licence, http://creativecommons.org/licenses/by/4.0/, as recorded in the arXiv licence metadata for this exact version and shown on the abstract page https://arxiv.org/abs/2510.06577v2. Full licence notice: \texttt{licenses/arxiv-2510.06577-CC-BY-4.0.txt}.\par\smallskip

\noindent\textit{Editorial scope.} Counted unit: the article's proposition estimate (source0.tex lines 585--591). The article's complete original proof of it is delivered with it (source0.tex lines 593--729, one \texttt{proof} environment of 137 lines). No mathematics is written, repaired or re-derived: the statement and the proof are the article's own lines, in the article's own notation, and no retained line is substituted or re-flowed.  Retained uncounted as the source-local context the proof needs: lines 234--236; lines 434--436; lines 442--444.  Nothing was omitted from any retained span, so the package carries no omission record.  No retained line contains a citation, so the wrapper carries no bibliography and no bibliography entry.  No retained line contains a figure environment, an image-inclusion command or drawing code, so no image is delivered and the package declares no asset dependency.  Editorial formatting only, in addition: an AMS \texttt{amsart} preamble that restates the article's own declarations and macros used by the retained lines, namely the environments \texttt{proposition} on separate counters, together with the standard packages those lines need.  The retained environments print in this document as Proposition 1 in order of appearance, whereas the article prints the same items with the numbering of its own full document.  The complete native source of the article as distributed in the arXiv e-print for this exact version is bundled unchanged as \texttt{sources/186618/source0.tex}.
\begin{equation}\label{scalar curvature problem}
F\Big(\nabla^2u+\frac{1-t}{n-2}(\Delta u)g+\frac{2-t}{2}|\nabla u|^2g-du\otimes du-A^{t}_{g}\Big)=\psi e^{2u}\,.
\end{equation}

\begin{equation}\label{relation}
\overline{F}^{ij}=F^{ij}+\frac{1-t}{n-2}\mathcal{F}\delta_{ij}.
\end{equation}

\begin{equation}\label{vanishes1}
0=\nabla_{i}G=\frac{1}{\beta}u_{k}u_{ki}+ \phi'u_{i}\,, \quad i\in \{1,\cdots,n\}\,,
\end{equation}

\begin{proposition}\label{int hessian estimate}
Assume  $\lambda(-g^{-1}A_{g}^{t})\in \Gamma$ for some $t<1$. Let $u\in C^{4}(M)$ be a solution to  \eqref{scalar curvature problem}. Then we have
\begin{equation*} 
\sup_{M}|\nabla^2u|\leq C_{2}\,,
\end{equation*}
where $C_{2}$ is a constant depends on $p,$  $C_{0}$, $C_{1}$, and  $ \psi$.
\end{proposition}

\begin{proof}
For each $x\in M$, let $S(T_{x}M)$ denote the unit tangent bundle of $M$ at $x$. For a large constant $B$ to be specified later, we define the quantity
\begin{equation*}
G_{0}=\max_{x\in M}\max_{\xi\in S(T_{x}M)}\bigg\{\nabla^2 u(\xi,\xi)+\frac{B}{2}|\nabla u|^{2}\bigg\}\,.
\end{equation*}
Our goal is to estimate $G_0$. Suppose $G_0$ is attained at a point $x_0 \in M$ and a unit vector field $\xi_0 \in S(T_{x_0} M)$ such that
\[
G_{0}=\nabla^2 u(\xi_{0},\xi_{0})+\frac{B}{2}|\nabla u|^{2}(x_{0})\,.
\]
In a neighborhood of $x_0$, we may choose a smooth orthonormal frame $\{e_1, \cdots, e_n\}$ satisfying:
\begin{itemize}\itemsep0.5em
\item[(i)] $e_{1}(x_0)=\xi_{0}$;
\item[(ii)] $\nabla_{e_i}e_j=0$ at $x_0$, i.e. $\Gamma_{ij}^k(x_0)=0$ for all $i,j$ and $k$;
\item[(iii)] the symmetric matrix $\{\nabla^2u(e_{i},e_{j})\}$ is diagonal at $x_0$.
\end{itemize}

Now consider the auxiliary function
\[
G=\nabla^2u(e_1,e_1)+\frac{B}{2}|\nabla u|^2=u_{11}-\Gamma_{11}^{l}u_{l}+\frac{B}{2}|\nabla u|^2\,.
\]
By the definition of $G_0$, the function $G$ attains its maximum $G_0$ at $x_0$.
The maximum principle  at $x_0$ then implies that
\begin{equation}\label{vanishes 1}
\begin{split}
0=\nabla_{i}G &= u_{11i}-\Gamma_{11}^{l}u_{li}-(e_{i}\Gamma_{11}^{l})u_{l}+\frac{B}{2}(e_{i}g^{kl})u_{k}u_{l}+Bg^{kl}(e_{i}u_{k})u_{l}\\
&=u_{11i}-(e_{i}\Gamma_{11}^{l})u_{l}+Bu_{ik}u_{k}\,,
\end{split}
\end{equation}
and
\begin{equation}\label{max principle}
\begin{split}
0\geq \overline{F}^{ij} u_{11ij}&-2\overline{F}^{ij}(e_{i}\Gamma_{11}^{j})u_{jj}-\overline{F}^{ij}(e_{j}e_{i}\Gamma_{11}^{l})u_{l}+\frac{B}{2}\overline{F}^{ij}e_{j}e_{i}(g^{kl})u_{k}u_{l}\\&+B\overline{F}^{ij}u_{ijl}u_{l}+B\overline{F}^{ii}u_{ii}^{2}\,,
\end{split}
\end{equation}
where ${\overline{F}^{ij}}$ is defined as in \eqref{relation}.

We first proceed to compute the last two terms in \eqref{max principle}. Differentiating the equation \eqref{scalar curvature problem} in the direction $e_l$ and using the fact that $\Gamma_{ij}^{k}=0$ at $x_0$, we deduce from \eqref{vanishes1} that
\begin{equation}\label{mix term}
\begin{split}
(\psi_{l}+2\psi u_{l})e^{2u}&=F^{ij}\big[(W[u])_{ij}\big]_{l}\\
&=F^{ij}\big(u_{ijl}-u_{k}(\Gamma_{ij}^{m})_{l} \big)-F^{ij}\big(u_{il}u_{j}+u_{i}u_{jl}+\nabla_{l}A^{t}_{ij}\big)\\
& \quad +\Big(\frac{1-t}{n-2}(u_{kkl}-u_{i}(\Gamma_{kk}^{i})_{l})+(2-t)u_{k} u_{kl}\Big)\mathcal{F}\\
&=\overline{F}^{ij}u_{ijl}-F^{ij}\big(2u_{il}u_{j}+\nabla_{l}A^{t}_{ij}+u_{k}(\Gamma_{ij}^{m})_{l}\big)\\
&\quad +\big((2-t)u_{k} u_{kl}-\frac{1-t}{n-2}u_{i}(\Gamma_{kk}^{i})_{l}\big)\mathcal{F}\,.
\end{split}
\end{equation}
Let us denote $W_{ij}=(W[u])_{ij}$ for simplicity. Then
\[
W_{ij}=u_{ij}+\frac{1-t}{n-2}(\Delta u)\delta_{ij}+\frac{2-t}{2}|\nabla u|^2\delta_{ij}-u_{i}u_{j}-A^{t}_{ij}\,.
\]
Since $W[u] \in \Gamma$ and by gradient estimates, we have
\[
0\leq \sum_{i=1}^{n}W_{ii}\leq \big(1+\frac{n(1-t)}{n-2}\big)\big(\Delta u-nC'\big)
\]
for some uniform positive constant $C'$. Therefore, we obtain
\[
\sum_{i=1}^{n}(u_{ii}-C')\geq 0\,.
\]
Note that $u_{11}$ is the largest eigenvalue of $\nabla^2u$, so for all $i=1,2,\cdots,n$,
\[
\begin{split}
|u_{ii}|&=|u_{ii}-C'+C'|\leq |u_{ii}-C'|+C'\\
&\leq  (n-1)(u_{11}-C')+C'=(n-1)u_{11}+nC'\,.
\end{split}
\]
It follows from \eqref{mix term} and the gradient estimates that
\begin{equation}
B\overline{F}^{ij}u_{ijl}u_{l}\geq -BC' (1+\mathcal{F}+u_{11}\mathcal{F} )\,.
\end{equation}

For the last term in \eqref{max principle}, we obtain
\begin{equation}\label{square}
\overline{F}^{ii}u_{ii}^{2}=\sum_{i=1}^{n}\Big( F^{ii}+\frac{1-t}{n-2}\mathcal{F}\Big)u_{ii}^{2}\geq \frac{1-t}{n-2}\mathcal{F} u_{11}^{2}
\end{equation}
For the second to fourth terms in \eqref{max principle}, we have
\begin{equation}\label{234}
-2\overline{F}^{ij}(e_{i}\Gamma_{11}^{j})u_{jj}-\overline{F}^{ij}(e_{j}e_{i}\Gamma_{11}^{l})u_{l} +\frac{B}{2}\overline{F}^{ij}e_{j}e_{i}(g^{kl})u_{k}u_{l}\geq - C'(1+u_{11}\mathcal{F}+B \mathcal{F})\,.
\end{equation}
Substituting \eqref{mix term}, \eqref{square}, and \eqref{234} into \eqref{max principle}, we deduce that
\begin{equation}\label{max principle v2}
\begin{split}
BC' (1+\mathcal{F}+u_{11}\mathcal{F} )\geq &\quad \overline{F}^{ij} u_{11ij}+\frac{B(1-t)}{n-2}\mathcal{F} u_{11}^{2}\,.
\end{split}
\end{equation}


Differentiating the equation \eqref{scalar curvature problem} in the $e_1$-direction twice, we obtain
\begin{equation*}
F^{ij,kl}e_{1}(W_{ij})e_{1}(W_{kl})+F^{ij}e_{1}e_{1}(W_{ij})=(\psi_{11}+4\psi_{1}u_{1}+4fu_{1}^2+2\psi u_{11})e^{2u}\,.
\end{equation*}
Since $F$ is concave on $\Gamma$, the first term on the left-hand side must be non-positive. Therefore, we have
\begin{equation}\label{fourth term}
F^{ij}e_{1}e_{1}(W_{ij})\geq -C'(1+u_{11})\,.
\end{equation}
A direct computation at $x_0$ gives
\[
\begin{split}
e_{1}e_{1}(W_{ij})=&-e_{1}e_{1}(A^{t}_{ij})+u_{ij11}+2u_{k1}e_{1}(\Gamma_{ij}^{k})+u_{k}e_{1}e_{1}(\Gamma_{ij}^{k})  \\
&+\frac{1-t}{n-2}\big(u_{kk11}g_{ij}+u_{kk}e_{1}e_{1}(g_{ij})-2u_{l1}e_{1}(\Gamma_{kk}^{l})+u_{r}e_{1}e_{1}(\Gamma_{kk}^{l}) \big)\\
&-(u_{i11}u_{j}+u_{i}u_{j11}+2u_{i1}u_{j1})+\frac{2-t}{2}e_{1}e_{1}(g^{kl})u_{k}u_{l}\delta_{ij}\\
&+(2-t)\big(u_{k11}u_{k}+u_{11}^{2}\delta_{ij}+|\nabla u|^2e_{1}e_{1}(g_{ij}) \big)\,.
\end{split}
\]
Substituting this into \eqref{fourth term} and using \eqref{vanishes 1} to replace terms involving $u_{11i}$, we obtain
\begin{equation}\label{fourth term v2}
-C'(1+u_{11})\leq \overline{F}^{ij} u_{11ij}+C'(1+u_{11})\mathcal{F}+(2-t)\mathcal{F}u_{11}^{2}\,.
\end{equation}
Combining \eqref{max principle v2} and \eqref{fourth term v2}, we deduce that
\begin{equation*}
\begin{split}
 \Big(\frac{B(1-t)}{n-2}-(2-t)\Big)\mathcal{F} u_{11}^{2}\leq &BC' (1+\mathcal{F}+u_{11}\mathcal{F} )+C'(1+u_{11})+C'(1+u_{11})\mathcal{F}\\
 =& (1+B)C'(1+u_{11})\mathcal{F}+C'(1+u_{11}+B)\,.
\end{split}
\end{equation*}

Since we have assumed $t<1$,  we may now choose $B$ to be sufficiently large such that
\[
\frac{B(1-t)}{n-2}-(2-t)=2.
\]
Then it follows that
\[
\big(2u_{11}^{2}- (1+B)C'(1+u_{11}) \big)\mathcal{F}\leq C'(1+u_{11}+B)\,.
\]
Without loss of generality, we may assume $u_{11}$ is sufficiently large so that 
\[
u_{11}^{2}- (1+B)C'(1+u_{11})\geq 0\,,
\]
since otherwise the desired bound already holds. Therefore, we infer that
\[
C'(1+u_{11}+B)\geq \mathcal{F}u_{11}^{2}\geq \mathscr{T}_{0}u_{11}^{2},
\]
where we have used the fact that $\mathcal{F} \geq \mathscr{T}_{0}$ for some positive constant $\mathscr{T}_{0}$. Indeed, by the Cauchy-Schwarz inequality, we have
\[
\mathcal{F}=\sum_{i=1}^{n}F^{ii}\geq n\Big(\prod_{i=1}^{f}F^{ii} \Big)^{1/n}\geq n\mathscr{T}^{1/n}=:\mathscr{T}_{0}.
\]
This yields the desired upper bound for $u_{11}(x_0)$, and hence an upper bound for $G_0$. The proof is then complete.
\end{proof}

\end{document}
